% Generated by publication/build_sections.py; do not edit by hand.
\paragraph{Informal verdict.}
\textbf{A --- Complete and apparently correct.} Confidence: \textbf{high (0.97)}.

\paragraph{Lean coverage.}
Q896.q896 formalizes the oriented plane as R x R with O translated to zero; represents each of the n indexed sets as a nonzero closed ray and the zones as connected components of the actual complement; retains ray distinctness and the no-open-half-plane hypothesis; takes the printed integer-periodic vertex family, global determinant inequalities modulo p, and exact M edge-vector equation; and proves p <= n for every positive n,p. For p>=3 it proves nonzero/distinct edges and disjoint open turning cones by a strict linear-functional minimum argument, forces a ray direction into every cone through zone closures, and obtains an injection Fin p -> Fin n. It separately covers p=1 and p=2. Equality/sharpness is optional; Example.lean gives a nonvacuous n=p=3 instance.

\subsection{Audited informal answer}
Q896 --- complete solution
Verified restatement and conventions
The canonical PDF states the following problem on physical page 23, in the Geometry section. There are no subparts, footnotes, or additional editorial remarks.
From a point \(O\), draw \(n\) distinct rays in an oriented affine plane. Their complement has \(n\) connected components \(Z_1,\ldots,Z_n\), called zones, and no zone is assumed to contain an open half-plane. A direct convex polygonal line with \(p\) sides is a \(p\)-periodic sequence \((A_t)_{t\in\mathbb Z}\) such that every vertex other than the endpoints of the oriented side \(A_kA_{k+1}\) lies strictly to its left:

\[
\det\!\left(\overrightarrow{A_kA_{k+1}},
             \overrightarrow{A_kA_\ell}\right)>0
\]

whenever \(\ell\not\equiv k,k+1\pmod p\). Define \(M_k\) by

\[
\overrightarrow{OM_k}=\overrightarrow{A_kA_{k+1}}.
\]

Assume that no zone \(Z_j\) has both \(M_k\) and \(M_{k+1}\) in its closure. Prove that \(p\le n\). Luc Pommeret
Here and below:


all polygon indices are read modulo \(p\);


``\(M\) is adherent to \(Z_j\)'' means \(M\in\overline{Z_j}\);


``period \(p\)'' need not initially mean least period, although the determinant condition forces it to be the least period when \(p\ge3\).


There is no mathematical typo in the statement. The only implicit convention is the modular interpretation of \(k+1\).

Final theorem
Let \(u_k\) be the direction of the \(k\)-th side:

\[
e_k:=\overrightarrow{A_kA_{k+1}},\qquad
u_k:=\frac{e_k}{\lVert e_k\rVert}\in S^1.
\]

Then the side directions \(u_0,\ldots,u_{p-1}\) occur exactly once around the unit circle in positive order. Moreover, the open positively oriented circular arc from \(u_k\) to \(u_{k+1}\) contains at least one of the \(n\) prescribed ray directions. These \(p\) arcs are disjoint, hence

\[
\boxed{p\le n}.
\]

In fact, more is true:


If \(p=n\), the side directions and the prescribed ray directions strictly alternate around the circle.


For every prescribed system of rays satisfying the half-plane hypothesis, an admissible polygon with \(p=n\) exists. Thus, for each fixed ray system, the maximum possible value of \(p\) is exactly \(n\).


The assumption that no zone contains an open half-plane is not needed for the upper bound itself.



Proof of the upper bound
We first treat the usual polygonal case \(p\ge3\). If the definition is read literally for \(p=1\) or \(p=2\), the original half-plane hypothesis implies \(n\ge3\), so the conclusion is immediate.
\begin{enumerate}
\item The determinant condition really gives a strictly convex polygon
\end{enumerate}
Set

\[
e_k=A_{k+1}-A_k.
\]

First, \(e_k\neq0\). Indeed, since \(p\ge3\), one may choose an index \(\ell\not\equiv k,k+1\), and then

\[
\det(e_k,A_\ell-A_k)>0
\]

would be impossible if \(e_k=0\).
The points \(A_0,\ldots,A_{p-1}\) are pairwise distinct. For if \(A_r=A_s\) with \(r\not\equiv s\), then either \(s\equiv r+1\), contradicting \(e_r\neq0\), or else

\[
\det(e_r,A_s-A_r)=0,
\]

contradicting the defining inequality.
Let

\[
K=\operatorname{conv}\{A_0,\ldots,A_{p-1}\}.
\]

For each \(k\), define the affine functional

\[
\Phi_k(X)=\det(e_k,X-A_k).
\]

The defining condition says

\[
\Phi_k(A_k)=\Phi_k(A_{k+1})=0,
\qquad
\Phi_k(A_\ell)>0
\quad
(\ell\not\equiv k,k+1).
\]

Consequently \(\Phi_k\ge0\) on \(K\). Moreover,

\[
K\cap\{\Phi_k=0\}=[A_k,A_{k+1}].
\]

Indeed, if \(X=\sum_\ell\lambda_\ell A_\ell\in K\) and \(\Phi_k(X)=0\), then

\[
0=\sum_\ell\lambda_\ell\Phi_k(A_\ell),
\]

where every summand is nonnegative and all summands except those for \(\ell=k,k+1\) are strictly positive whenever their coefficient is nonzero. Hence \(X\) is a convex combination of \(A_k\) and \(A_{k+1}\).
Thus every segment \([A_k,A_{k+1}]\) is an exposed edge of \(K\). Its endpoints are vertices of \(K\), so \(K\) has precisely the \(p\) vertices \(A_0,\ldots,A_{p-1}\). A convex polygon with \(p\) vertices has \(p\) edges, and we have already exhibited \(p\) distinct edges. Therefore

\[
[A_0,A_1], [A_1,A_2],\ldots,[A_{p-1},A_0]
\]

are exactly the boundary edges of \(K\).
Since \(K\) lies to the left of every directed side \(A_kA_{k+1}\), the sequence traverses the boundary of \(K\) in the positive, counterclockwise direction.
In particular, the original condition has ruled out any smaller period.

\begin{enumerate}
\item The side directions make exactly one positive revolution
\end{enumerate}
For every \(k\),

\[
\det(e_k,e_{k+1})>0.
\]

Indeed, applying the defining inequality for the \(k\)-th side to \(A_{k+2}\) gives

\[
\begin{aligned}
0
&<
\det(e_k,A_{k+2}-A_k)\\
&=
\det(e_k,e_k+e_{k+1})\\
&=
\det(e_k,e_{k+1}).
\end{aligned}
\]

Hence the positive turning angle \(\delta_k\) from \(e_k\) to \(e_{k+1}\) belongs to \((0,\pi)\).
Let \(\alpha_{k+1}\in(0,\pi)\) be the interior angle of \(K\) at \(A_{k+1}\). Then

\[
\delta_k=\pi-\alpha_{k+1}.
\]

Triangulating the convex \(p\)-gon from one vertex gives

\[
\sum_{k=0}^{p-1}\alpha_{k+1}=(p-2)\pi.
\]

Therefore

\[
\sum_{k=0}^{p-1}\delta_k
=p\pi-(p-2)\pi
=2\pi.
\]

We may consequently choose real arguments

\[
\theta_0<\theta_1<\cdots<\theta_p=\theta_0+2\pi
\]

such that

\[
u_k=(\cos\theta_k,\sin\theta_k),
\qquad
\theta_{k+1}-\theta_k=\delta_k
\quad(0\le k<p),
\]

where \(u_p=u_0\).
For each \(k\), let

\[
I_k
=
\bigl\{(\cos\theta,\sin\theta):
       \theta_k<\theta<\theta_{k+1}\bigr\}.
\]

Thus \(I_k\) is the open positively oriented arc from the direction of side \(k\) to the direction of side \(k+1\).
The arcs

\[
I_0,\ldots,I_{p-1}
\]

are pairwise disjoint, and together they cover the unit circle except for the \(p\) side directions \(u_0,\ldots,u_{p-1}\).

\begin{enumerate}
\item Every turning arc contains a prescribed ray
\end{enumerate}
Let \(S^1\) be the unit circle centered at \(O\). Each of the \(n\) prescribed rays meets \(S^1\) in exactly one point. Denote the set of these direction points by

\[
B=\{b_1,\ldots,b_n\}\subset S^1.
\]

The connected components of \(S^1\setminus B\) correspond exactly to the planar zones \(Z_j\). If \(C_j=Z_j\cap S^1\), then

\[
\overline{Z_j}
=
\{O\}\cup
\{r u:r>0,\ u\in\overline{C_j}\}.
\]

Fix \(k\). Suppose, seeking a contradiction, that

\[
I_k\cap B=\varnothing.
\]

The arc \(I_k\) is connected and is contained in \(S^1\setminus B\), so it lies in one connected component \(C_j\). Its two endpoints \(u_k\) and \(u_{k+1}\) therefore belong to \(\overline{C_j}\).
But

\[
M_k=\lVert e_k\rVert u_k,
\qquad
M_{k+1}=\lVert e_{k+1}\rVert u_{k+1},
\]

so both \(M_k\) and \(M_{k+1}\) lie in \(\overline{Z_j}\). This contradicts the hypothesis.
We have proved

\[
I_k\cap B\neq\varnothing
\qquad\text{for every }k.
\]

Since the \(p\) arcs \(I_k\) are pairwise disjoint, the ray direction chosen from each \(I_k\) must be different. Hence \(B\) has at least \(p\) elements:

\[
n=\lvert B\rvert\ge p.
\]

Therefore

\[
\boxed{p\le n}.
\]


Equality case
Suppose \(p=n\). Each of the \(p\) disjoint arcs \(I_k\) contains at least one of the \(n\) ray directions. Consequently:


each \(I_k\) contains exactly one ray direction;


no ray direction coincides with a side direction \(u_k\).


After cyclic relabeling, the directions therefore occur in the strict order

\[
u_0,\ b_0,\ u_1,\ b_1,\ldots,
u_{p-1},\ b_{p-1},\ u_0.
\]

Thus:

Equality characterization.
If \(p=n\), the side directions of the polygon and the prescribed rays strictly alternate around \(O\). Equivalently, every zone contains exactly one side direction.

Under the original half-plane hypothesis, the converse reformulation is also exact. Every zone has angular width \(<\pi\), and every polygonal turning angle \(\delta_k\) is \(<\pi\). Hence \(M_k\) and \(M_{k+1}\) belong to the closure of one zone exactly when the open turning arc \(I_k\) contains no prescribed ray.

Sharpness for every prescribed ray configuration
We now prove that, for every fixed system of \(n\) rays satisfying the hypotheses, equality \(p=n\) can be achieved.
Order the ray arguments positively:

\[
\beta_0<\beta_1<\cdots<\beta_{n-1}<\beta_n:=\beta_0+2\pi.
\]

Put

\[
\gamma_i=\beta_{i+1}-\beta_i.
\]

The zone between the rays of arguments \(\beta_i\) and \(\beta_{i+1}\) has angular width \(\gamma_i\). Since it contains no open half-plane,

\[
0<\gamma_i<\pi.
\]

Choose the midpoint direction of each zone:

\[
\theta_i=\frac{\beta_i+\beta_{i+1}}2,
\qquad
u_i=(\cos\theta_i,\sin\theta_i).
\]

Then

\[
\theta_{i+1}-\theta_i
=
\frac{\gamma_i+\gamma_{i+1}}2
\in(0,\pi),
\]

with cyclic indices. Thus the \(u_i\)'s occur positively around \(S^1\), and every gap between two consecutive \(u_i\)'s has length \(<\pi\).
Consequently the points \(u_i\) are not contained in any closed semicircle. It follows that

\[
0\in
\operatorname{int}
\operatorname{conv}\{u_0,\ldots,u_{n-1}\}.
\]

For completeness, if the origin were not in the interior of that convex hull, there would be a line through the origin supporting or separating the convex hull, placing all \(u_i\)'s in one closed half-plane, hence in a closed semicircle.
We next choose positive side lengths. Write

\[
s=\sum_{i=0}^{n-1}u_i.
\]

Because the origin is an interior point of the convex hull, for sufficiently small \(t>0\),

\[
-ts\in\operatorname{conv}\{u_0,\ldots,u_{n-1}\}.
\]

Thus there are \(\mu_i\ge0\), with \(\sum_i\mu_i=1\), such that

\[
-ts=\sum_{i=0}^{n-1}\mu_i u_i.
\]

Set

\[
\ell_i=t+\mu_i>0.
\]

Then

\[
\sum_{i=0}^{n-1}\ell_i u_i
=
t\sum_i u_i+\sum_i\mu_i u_i
=ts-ts
=0.
\]

Define

\[
e_i=\ell_i u_i,\qquad
A_0=0,\qquad
A_{k+1}=A_k+e_k.
\]

The relation \(\sum e_i=0\) gives \(A_n=A_0\).
It remains to verify the full strict convexity condition, rather than merely local positive turns.
Fix \(k\), extend the \(\theta_i\)'s by

\[
\theta_{i+n}=\theta_i+2\pi,
\]

and let \(2\le m\le n-1\). Then

\[
A_{k+m}-A_k
=
\sum_{r=0}^{m-1}e_{k+r}.
\]

Consequently

\[
\frac{1}{\ell_k}
\det(e_k,A_{k+m}-A_k)
=
\sum_{r=1}^{m-1}
\ell_{k+r}
\sin(\theta_{k+r}-\theta_k).
\tag{1}
\]

If \(\theta_{k+m-1}-\theta_k\le\pi\), all terms on the right of \(1\) are nonnegative, and the first one is strictly positive because

\[
0<\theta_{k+1}-\theta_k<\pi.
\]

Hence \(1\) is positive.
If instead

\[
\theta_{k+m-1}-\theta_k>\pi,
\]

use \(\sum e_i=0\) to rewrite \(1\) as

\[
-\sum_{r=m}^{n-1}
\ell_{k+r}
\sin(\theta_{k+r}-\theta_k).
\]

Every angle in this last sum lies strictly between \(\pi\) and \(2\pi\), so every sine is negative. The displayed expression is therefore strictly positive.
Thus, for every \(\ell\not\equiv k,k+1\),

\[
\det(e_k,A_\ell-A_k)>0.
\]

The sequence \((A_k)\) is a direct strictly convex \(n\)-gon.
Finally, \(u_i\) is the midpoint direction of \(Z_i\), so

\[
M_i=e_i=\ell_i u_i\in Z_i.
\]

Hence \(M_i\) and \(M_{i+1}\) lie in the interiors of two different zones. An interior point of a zone belongs to the closure of no other zone, so no zone closure contains both.
We have constructed an admissible polygon with

\[
p=n.
\]

Therefore, for every prescribed ray system satisfying the statement's hypotheses,

\[
\boxed{\max p=n}.
\]


Degeneracies and hypothesis audit
The half-plane assumption is unnecessary for the upper bound
The counting proof only used the fact that the zones correspond to the connected components of the circle after removing the ray directions. It did not use their angular widths.
Even if \(p=2\) is allowed, the conclusion remains valid without the half-plane assumption. The two side directions are antipodal. If one of the two open semicircles between them contained no prescribed ray, the two side-vector points would belong to the closure of one zone. Hence each open semicircle contains a different ray and \(n\ge2=p\).
For \(p=1\), the only side vector is zero, so \(M_0=O\), which belongs to the closure of every zone; the separation hypothesis cannot hold.
Thus the half-plane condition can be deleted completely from the upper-bound theorem. Its useful role is to guarantee the universal equality construction above.
The separation condition is essential
Take three equally spaced rays and any direct convex quadrilateral. If the condition concerning consecutive \(M_k\)'s is removed, then \(p=4>3=n\). Thus some separation of consecutive side directions by the prescribed rays is indispensable.
Global convexity is essential
Merely requiring all local turns to be positive does not suffice.
Take three rays of arguments \(0,2\pi/3,4\pi/3\), and let

\[
u_0=
\left(\frac12,\frac{\sqrt3}{2}\right),\qquad
u_1=(-1,0),\qquad
u_2=
\left(\frac12,-\frac{\sqrt3}{2}\right).
\]

These are the midpoint directions of the three zones and satisfy

\[
u_0+u_1+u_2=0.
\]

Consider the six successive edge vectors

\[
u_0,\quad
2u_1,\quad
3u_2,\quad
4u_0,\quad
3u_1,\quad
2u_2.
\]

Their sum is

\[
5(u_0+u_1+u_2)=0,
\]

so they form a closed six-sided polygonal line. Every consecutive pair of directions belongs to two different zone interiors, and every local turn is \(2\pi/3>0\). Nevertheless \(p=6>3\).
The chain is not globally convex. For example, with \(A_0=0\),

\[
A_3=u_0+2u_1+3u_2=(0,-\sqrt3),
\]

and therefore

\[
\det(u_0,A_3-A_0)
=
-\frac{\sqrt3}{2}<0.
\]

This illustrates exactly what the full determinant hypothesis contributes: it forces the side directions to make one, and only one, revolution around the circle.

Conclusion
The strict convexity condition orders the \(p\) side directions once around \(S^1\). Between each two consecutive side directions, the hypothesis forces at least one of the prescribed ray directions. These \(p\) angular intervals are disjoint, so they require at least \(p\) distinct rays:

\[
\boxed{p\le n}.
\]

Moreover, equality is attainable for every admissible configuration of rays, and in the equality case the ray directions and side directions strictly alternate.
